\begin{textsample}{POS Dim 3 – go – Score 22.00 – go/go\_inf\_16.txt}  \label{ex:f3_pos_047}
3.2.2.
Organização da Ação Pedagógica No desenvolvimento das ações pedagógicas relacionadas aos sentidos, saberes e conhecimentos do campo “ Corpo, \textbf{gestos} e movimentos ” é importante ter clareza que o corpo das pessoas, com as quais a \textbf{criança} convive transmite muitas informações para ela : o tom de voz, as expressões faciais, o olhar, a postura, os \textbf{gestos} e os \textbf{cuidados} que dispensa a ele.
Isso quer dizer que a importância que o professor atribui ao seu corpo, afeta as \textbf{crianças} com as quais trabalha ( TRINDADE, 2017 ).
Os professores, em sua maioria, não se movimentam e buscam reprimir o movimento das \textbf{crianças}.
Mas, como possibilitar às \textbf{crianças} variados \textbf{gestos} e movimentos se os professores não os tiverem vivenciado?
Como perceber as sensações dos outros sem se conectar com as suas?
Como superar dificuldades e possibilidades expressas no corpo do professor?
É papel do professor se sentar no \textbf{chão}, \textbf{brincar} com as \textbf{crianças}, se deixar molhar, pintar e afagar pelas \textbf{crianças}, bem como afagá -las ; ter disponibilidade às experiências corporais ; usar \textbf{vestimentas} que favoreçam se movimentar junto com as \textbf{crianças} ; além de se permitir aprender na relação com elas.
A promoção de oportunidades variadas de exploração de diferentes ambientes é fundamental para a \textbf{criança} se apropriar de sua \textbf{movimentação}, explorando -a e recriando -a a partir da descoberta e ampliação dos modos de diálogo entre corpo e espaço ; das formas produzidas a partir das \textbf{movimentações} e coreografias, linhas retas ou diagonais, da ocupação e uso do espaço com o corpo e das composições espaciais que podem ser feitas a partir de cenários imaginários criados e recriados, tais como uma floresta, um espaço sideral.
Nesse processo é preciso considerar e aproveitar as realidades de vida, de culturas, dos contextos de cada instituição educacional, a fim de que as oportunidades sejam oferecidas à \textbf{criança}.
Considerar também que o corpo da \textbf{criança} é seu patrimônio e que os \textbf{adultos} são responsáveis por sua integridade.
Assim, conhecimento do grupo com o qual se trabalha, prudência e consciência dos perigos que o ambiente pode oferecer - pisos escorregadios, estado de conservação dos móveis etc. - são condições indispensáveis para realizar atividades de práticas corporais sem que isso ofereça riscos à integridade física da \textbf{criança}, já que os professores são responsáveis por sua segurança e bem-estar quando ela está na instituição.
A compreensão de que o movimento faz parte do ambiente da instituição de Educação \textbf{Infantil} é indispensável para superar a perspectiva de manter a \textbf{criança} em fila, quieta, aguardando o comando do \textbf{adulto} que perpetua a ideia de repressão do movimento de corpo inteiro, destituído de totalidade.
O tão imperativo silêncio dos corpos “ Fica quieta!
Não se mexe!
Senta! ”"Deita a cabeça na mesa!"precisa dar lugar ao movimento e a criatividade \textbf{infantil}.
A proposição de diferentes \textbf{brincadeiras}, em variados espaços, contribui para a superação dessas práticas.
Contudo, é preciso ter clareza de que a \textbf{criança} não nasce sabendo \textbf{brincar} ou jogar, sendo necessário criar contextos significativos na instituição educativa que possibilitem o desenvolvimento dessa ação pela \textbf{criança} na relação com os outros, de forma que amplie seu repertório de \textbf{brincadeiras}.
O trabalho com o movimento deve partir do patrimônio cultural e corporal do grupo e se iniciar com as \textbf{brincadeiras} que a \textbf{criança} conhece, aquelas cultivadas em seu ambiente familiar, mas não com o objetivo de repeti -las ou simplesmente passar o tempo e sim ampliar seu repertório de \textbf{brincadeiras}, \textbf{gestos} e movimentos.
Cabe aos professores de Educação \textbf{Infantil} \textbf{brincar} de coisas que a \textbf{criança} \textbf{brinca}, mas também ensinar outras \textbf{brincadeiras} e outras formas de \textbf{brincar}, promover \textbf{brincadeiras} em que \textbf{meninos} e \textbf{meninas} \textbf{brinquem} juntos, reconstruindo regras, conhecendo \textbf{brincadeiras} e danças do passado e do presente, cantigas de roda e \textbf{brincadeiras} de origem indígena, europeia, africana ; esportes com e sem materiais ; situações que simulem os deslocamentos e comportamentos exigidos no trânsito, entre outras possibilidades.
As \textbf{brincadeiras} têm valor cultural e devem ser vivenciadas plenamente com esse fim.
A \textbf{criança} \textbf{brinca} de várias maneiras em diferentes momentos, e em todas elas a qualidade da mediação do \textbf{adulto}, na organização de materiais e espaços, no planejamento dos tempos, nos desafios propostos, na organização do grupo, é um fator importante para o desenvolvimento \textbf{infantil}.
Outro aspecto a ser considerado na organização da ação pedagógica é a redução dos espaços para a \textbf{criança} se movimentar, que se dá por motivos, tais como : o tamanho das residências que é cada vez \textbf{menor}, muitas delas não têm quintal ; não há muito acesso a praças, parques ou áreas de lazer e ainda há o medo da violência, que impede que a \textbf{criança} se desloque pelos diferentes lugares do bairro, da cidade, da comunidade, levando ao confinamento \textbf{doméstico}.
Comumente se vê \textbf{crianças} em salas por longos períodos, tendo que ficar quietas e silenciosas, sendo necessário propor alternativas, abrindo mão do medo da mobilidade que tanto profissionais quanto famílias às vezes têm.
É importante que as \textbf{crianças} tenham possibilidades de fazer deslocamentos e movimentos amplos nos diferentes espaços da instituição.
Acontecem com frequência situações em que a \textbf{criança} apresenta ao professor perguntas e questionamentos sobre seu corpo e sobre sexualidade.
Mas, como tratar essas questões com a \textbf{criança}?
Primeiro, deve -se responder exatamente o que ela perguntou, sanando apenas as dúvidas questionadas e usando linguagem simples e clara.
É sempre oportuno ter uma resposta mais completa, caso prossiga com as perguntas.
Quando uma \textbf{criança} indagar alguma coisa sobre a genitália, por exemplo, dizer sempre o nome correto, não usar apelidos ou termos jocosos, não fazer julgamentos e não \textbf{demonstrar} espanto ou constrangimento.
Em geral, a \textbf{crianças} se contenta com as primeiras respostas.
As dúvidas e \textbf{curiosidades} podem ser sanadas trabalhando com bonecos, quando se pode nomear as partes do corpo, informar se é masculino ou feminino, refletir sobre o crescimento e as mudanças corporais.
Reiterando, é fundamental responder todas as questões com simplicidade, mas corretamente.
Se o professor fica constrangido nessas situações, pode recorrer a outro profissional da instituição para encontrar meios e responder as perguntas das \textbf{crianças}, que não podem ficar sem resposta, evitando informações estereotipadas, preconceituosas e incorretas.
Na medida em que essas dúvidas e questionamentos vão surgindo, é preciso estabelecer um diálogo aberto e respeitoso com as famílias, pois ao oferecer respostas aos questionamentos das \textbf{crianças}, é preciso que as famílias estejam a par disso e saibam que será um diálogo sem preconceito ou discriminação de qualquer natureza, que visa atender as \textbf{curiosidades} das \textbf{crianças}, considerando que a percepção que os \textbf{adultos} têm sobre essas questões, não é a mesma das \textbf{crianças}.
As instituições também têm o papel de ajudar a \textbf{criança} a diferenciar uma aproximação saudável, respeitosa, fruto de afetividade e \textbf{cuidado}, daquela que pode se constituir em situações de risco ou de abuso.
É preciso orientar a \textbf{criança} quanto ao contato físico com \textbf{adultos} ou com outras \textbf{crianças}, auxiliando -as a diferenciar \textbf{gestos} e toques de carinho e \textbf{cuidado}, de situações abusivas.
Orientá -la que ninguém deve tocar suas partes íntimas e nem elas devem tocar as partes íntimas de outra pessoa e que, caso isso aconteça, deve pedir \textbf{ajuda} a um \textbf{adulto} de confiança.
A \textbf{criança} também precisa ter a oportunidade de tornar -se o principal agente nos seus \textbf{cuidados} pessoais.
A instituição educativa, \textbf{progressivamente}, pode levar a \textbf{criança} a adquirir consciência e autonomia nos \textbf{cuidados} com seu próprio corpo, oportunizando que ela vivencie a \textbf{rotina} com sentido e significado, organizada e desenvolvida a partir da mediação e \textbf{acompanhamento} dos \textbf{adultos}, levando -a a participar dos \textbf{cuidados} com seu corpo, fazer escolhas, tornando -se \textbf{progressivamente} responsável e autônoma no que diz respeito a sua \textbf{higiene}, seus \textbf{hábitos} alimentares, seu bem estar e autocuidado.
Noutro sentido, Greguol ( 2017 ) aponta que a \textbf{criança} com deficiências é privada de vivenciar experiências motoras e recreativas.
Alerta que a falta de oportunidades de sociabilizar -se com outras \textbf{crianças}, de \textbf{brincar} e de frequentar ambientes estimulantes pode causar perdas sensíveis e, inclusive, irreversíveis na função motora.
É necessária a mudança do foco, enfatizando -se o que a \textbf{criança} pode realizar e não as incapacidades, fomentando suas potencialidades, o que pode resultar em proposições criativas que permitam a estimulação do repertório motor e o desenvolvimento corporal de todas as \textbf{crianças}.
Assim, é importante que \textbf{crianças} com deficiências vivam plenamente a \textbf{rotina} da instituição, participando de situações de \textbf{movimentação}, de \textbf{brincadeiras}, de passeios ; experimentando o movimento dançado ; saindo de suas salas ; se relacionando e interagindo com as \textbf{crianças} de outras turmas ou agrupamentos ; tendo acesso aos \textbf{brinquedos}, aos livros etc.
O que se deve é possibilitar que as \textbf{crianças}, todas elas, se movimentem de modo a experimentar o mundo de corpo inteiro.

% matched lemmas: acompanhamento, adulto, ajuda, brincadeira, brincar, brinquedo, chão, criança, cuidado, curiosidade, demonstrar, doméstico, gesto, higiene, hábito, infantil, menino, movimentação, pequeno, progressivamente, rotina, vestimenta
\end{textsample}
